\chapter{Introduzione}

\section{Contesto e scenario generale}
\subsection{Hypervisor}
In ambito infrastrutturale IT, una macchina di tipo server opera grazie all'impiego di sistemi operativi specialistici, noti come \textit{Hypervisor} \cite{IBM_hypervisor_overview}, specializzati nella virtualizzazione multipla di ambienti disaccoppiati tra loro e su cui gireranno i vari servizi.


\subsection{Model Context Protocol (MCP)}
Un qualsiasi \textit{LLM} può interagire con software classici interpretandoli come tools e instaurando con essi una comunicazione simil API, ma molto più stabdardizzata.
In questo scenario l'azienda \textit{Anthropic} ha sviluppato un protocollo chiamato \textit{MCP} \cite{mcp_anthropic_announcement}, che permette interazioni strutturate basate sulla logica di un sistema client-server, rispettivamente \textit{LLM} e il layer \textit{MCP}.

Attualmente il progetto \textit{MCP} è opensource ed è stato donato da \textit{Anthropic}, per garantirne l'indipendenza, alla \textit{Agentic AI Foundation} \cite{mcp_agentic_ai_foundation} \cite{mcp_project_page}, un fondo gestito dalla \textit{Linux Foundation}.

\section{Problematiche che si propone di risolvere}
La gestione di un \textit{hypervisor} richiede competenze sistemistiche verticali e l'uso di interfacce GUI o CLI spesso complesse e poco user-friendly.
In questo lavoro di tesi si propone di realizzare e valutare un \textit{MCP server} con il ruolo di bridge architetturale tra un \textit{LLM} e un hypervisor aggiungendo così una nuova interfaccia per l'utente basata sul linguaggio naturale.

\section{Obiettivi della tesi}
L'architettura non si limita ad interfacciare un \textit{LLM} ad un hypervisor, ma prevede soluzioni volte a irrobustire questa interfaccia con scelte di design che gestiscono operazioni potenzialmente distruttive e mantengono il controllo sulle azioni dell'\textit{LLM}.
